\documentclass{article}%
\usepackage[T1]{fontenc}%
\usepackage[utf8]{inputenc}%
\usepackage{lmodern}%
\usepackage{textcomp}%
\usepackage{lastpage}%
\usepackage{geometry}%
\geometry{margin=2.2cm}%
%
\usepackage{amsmath}%
\usepackage{amssymb}%
\usepackage{booktabs}%
\usepackage{longtable}%
\usepackage{array}%
\usepackage{hyperref}%
\hypersetup{hidelinks}%
\title{ITA 2026 {-} Física (claude\_opus\_5\_5) x Poliedro\\[4pt]\large Comparação com o gabarito Poliedro Resolve\thanks{Gerado com gpt-resolve: \url{https://github.com/lgabs/gpt-resolve}}}%
\date{10/11/2026}%
%
\begin{document}%
\normalsize%
\maketitle%
\section*{Resumo}\textbf{Nota média: 9,0/10} (20 de 22 itens com a resolução correta; cada item vale o mesmo dentro da questão).\par\medskip%
\begin{longtable}{@{}c c p{0.78\textwidth}@{}}\toprule Questão & Nota & Resumo \\ \midrule
1 & 10,0 & A resolução e o gabarito do Poliedro concordam nos três itens. Em (a), v$^2$(1 - $\mu$ tan$\theta$)/{[}g(tan$\theta$ + $\mu$){]} é a mesma expressão que (v$^2$/g)(cos$\theta$ - $\mu$ sen$\theta$)/(sen$\theta$ + $\mu$ cos$\theta$), e as duas equivalem a v$^2$/{[}g tan($\theta$ + arctg $\mu$){]}. Em (b), v$^2$/(g tan$\theta$) é o mesmo que (v$^2$/g)cotg $\theta$. Em (c), ambos dão v$^2$/($\mu$g). Os dois usam o mesmo modelo físico: a normal fica inclinada de $\theta$ em relação à vertical, rumo ao centro do trem, e no limite a roda interna tem N = 0 enquanto a externa está na iminência de escorregar. A resolução é mais completa: justifica pelo cone de atrito que esse é o caso extremo e indica que a resposta só vale para $\mu$ tan$\theta$ < 1. Ficam apenas pontos menores: a leitura de 'perder o contato' no item (c) e a hipótese implícita de desprezar os torques. \\
2 & 5,0 & Item a: a resolução e o gabarito coincidem em forma equivalente, já que M(L/2-d)/(L-2d) = M/2, e conferi numericamente que as duas normais estão corretas. Item b: os dois dão a mesma soma dos atritos, f\_A + f\_B = (M+m+m\_C)g sen$\theta$, rampa acima. Para cada atrito separado, o gabarito conclui que o problema é indeterminado, e isso está fisicamente correto: os atritos estão sobre a mesma linha de ação e não entram na equação de momentos. A resolução também diz que só a soma é determinada, mas depois apresenta em destaque f\_i = N\_i tan$\theta$, que vale só se os dois apoios tiverem o mesmo $\mu$ e a barra estiver na iminência de deslizar ($\mu$ = tan$\theta$). O enunciado não dá nenhuma das duas condições. Por isso o gabarito está mais correto no item b, e a resolução deve receber crédito parcial pela soma. Uma ressalva: o enunciado pede o atrito \textbackslash{}"em cada apoio\textbackslash{}" sem dar $\mu$, então a banca pode ter tido em mente uma divisão específica, e isso vale levar ao professor. \\
3 & 10,0 & A resolução e o gabarito Poliedro concordam nos dois itens. No item a, ambos dão a componente radial M{[}g cos $\theta$ + $\omega$$^2$(L-2a)/2{]}, apontando para o pivô, ou seja, para o lado do trecho curto. No item b, ambos dão $\theta$ = arccos{[}-$\omega$$^2$(L-2a)/(2g){]}; a resolução só escreve o argumento como $\omega$$^2$(2a-L)/(2g). Refiz a conta pelo teorema do CM e pela integral da força centrífuga no referencial girante, e a conferência numérica bate. A resolução ainda acrescenta a condição de existência $\omega$$^2$(L-2a) $\leq$ 2g. O único reparo é de redação: ela chama de 'módulo' uma expressão que pode ficar negativa para $\theta$ > 90$^\circ$. Os dois itens estão corretos. \\
4 & 10,0 & A resolução está certa nos dois itens e bate com o gabarito do Poliedro. No item a, as duas dão y$_{2}$ = 2b ou y$_{2}$ = -4b. No item b, as duas põem a estrela no foco (c,0) com y$_{2}$ = 2b, e os termos da estrela coincidem: -(GMm/5){[}16/$\surd$((3a-4c)$^2$+b$^2$) + 1/$\surd$((a-c)$^2$+4b$^2$){]} é o mesmo que -GMm$_{2}${[}4/d$_{1}$ + 1/d$_{2}${]}. A resolução acrescenta a interação m$_{1}$--m$_{2}$ e a declara desprezível; com isso fica mais completa, e o resultado é equivalente ao do gabarito. O ponto que merece atenção é o próprio enunciado: a premissa de que o CM continua na órbita original é fisicamente inconsistente com a separação dada, e o Poliedro sugere anular a questão. \\
5 & 5,0 & Conferi o enunciado no PDF original (p. 6). No item a, resolução e gabarito são equivalentes: W = (M1 - M2)g$\alpha$L$\Delta$T = 70(M1 - M2)g$\alpha$L; a resolução apenas substituiu g = 10 (700$\alpha$L(M1 - M2)). No item b, o gabarito está mais correto. Ele inclui no calor absorvido o termo mg$\alpha$L/2, que corresponde à subida do centro de massa da própria viga, e a resolução omite esse termo. A diferença numérica é desprezível (cerca de 10$^{-}$$^{8}$ relativo) e a forma aproximada da resolução coincide com o limite do gabarito. Mesmo assim, a expressão completa da resolução é incompleta, o que pode custar uma pequena parte da nota do item b. \\
6 & 10,0 & A resolução está certa e concorda com o gabarito Poliedro nos três itens. a) L = 85 cm. b) G = -16 (módulo 16), com imagem invertida. c) A imagem final é real e invertida, fica a 165/32 $\approx$ 5,2 cm depois da ocular e é 16 vezes menor que o objeto (A = -1/16 $\approx$ -0,063). Refiz as contas com frações exatas e os valores conferem. A única diferença é o arredondamento do gabarito (-0,063 contra -0,0625 exato). A resolução também mostra que o aumento lateral de um sistema afocal é -f2/f1, o que confirma o item (c) por outro caminho. \\
7 & 10,0 & Na questão 7, a resolução e o gabarito do Poliedro dão a mesma resposta nos dois itens: a) I = I0$\cdot$x\^{}(N+1)$\cdot$cos\^{}(2N)($\varphi$/N); b) r = sqrt(P/($\pi$$\cdot$I0$\cdot$x\^{}(N+1)$\cdot$cos\^{}(2N)($\varphi$/N))). Refiz as contas: o 1º polarizador só aplica a perda x, e cada um dos N seguintes aplica x$\cdot$cos$^2$($\varphi$/N). Depois, P = I$\cdot$$\pi$r$^2$ dá o raio. Os dois estão certos. Há apenas observações menores: a leitura de como x se aplica, a hipótese de que o analisador só mede potência e a frase imprecisa sobre conservação de energia na resolução. \\
8 & 10,0 & A resolução e o gabarito Poliedro Resolve concordam nos dois itens: a) P\_min = 35 W; b) $\Delta$t\_max = 2,24$\pi$$^2$$\cdot$10$^{-}$$^{6}$ s = (56/25)$\pi$$^2$$\cdot$10$^{-}$$^{6}$ s $\approx$ 2,2$\cdot$10$^{-}$$^{5}$ s. Conferi as contas em Python e o enunciado contra o PDF. A resolução justifica o 'máximo' com mais rigor (M $\leq$ L no acoplamento perfeito) e verifica a coerência com o item a (5 V $\times$ 7 A = 35 W). Não há divergência de resultado. \\
9 & 10,0 & A resolução está certa nos dois itens e confere com o gabarito Poliedro. a) d = 600$\surd$2 nm $\approx$ 848,5 nm, porque $\theta$ = 45$^\circ$ para o laranja. b) x = 20$\surd$14/7 cm $\approx$ 10,7 cm para o violeta, com sen $\theta$' = $\surd$2/3. Refiz as contas em Python e os valores batem. A única diferença é o caminho: a resolução usa tg $\theta$ e o gabarito usa sen $\theta$ = y/$\surd$(D$^2$+y$^2$), o que é equivalente. Os arredondamentos (849 nm $\approx$ 8,5$\times$10$^{-}$$^{7}$ m) não mudam nada. Também é correto o alerta da resolução de que a aproximação de ângulos pequenos daria 13,3 cm, valor errado. \\
10 & 10,0 & Resolução e gabarito coincidem nos dois itens e ambos estão certos. Item a: 12,09375 anos $\approx$ 12,1 anos. Item b: tempo próprio da nave 2 igual a 9,675 anos $\approx$ 9,7 anos. Refiz as contas com frações exatas em python e também confirmei o item b pela transformação de Lorentz. A transcrição confere com o original (página 11). Só há ressalvas de interpretação de baixa importância. \\
\bottomrule\end{longtable}%
\section*{Principais pontos de dúvida}%
\begin{itemize}%
\item%
\textbf{Q2} \textit{(statement)} O item b pede para "calcular a força de atrito em cada apoio" e lista os parâmetros permitidos sem coeficiente de atrito. Isso sugere que a banca esperava valores individuais, talvez proporcionais às normais (f\_i = N\_i tan$\theta$, como na resolução) ou divididos igualmente. Pela estática de corpo rígido, porém, só a soma é determinada, como diz o Poliedro. Vale levar ao professor qual critério a banca do ITA deve usar. Se o gabarito oficial der f\_i = N\_i tan$\theta$, a resolução fica com a resposta esperada.%
\item%
\textbf{Q2} \textit{(solution)} No item b, a resolução põe em destaque f\_A = N\_A tan$\theta$ e f\_B = N\_B tan$\theta$ com base na hipótese de deslizamento iminente ($\mu$ = tan$\theta$) e $\mu$ igual nos dois apoios. Nenhuma das duas vem do enunciado. A ressalva aparece no texto, mas a resposta final dá a impressão de um resultado determinado. Pelo critério do gabarito Poliedro (indeterminado), esses valores não se justificam.%
\item%
\textbf{Q4} \textit{(statement)} A questão depende de uma premissa que o enunciado não declara: o CM de m$_{1}$+m$_{2}$ continua sobre a elipse original. Com a separação dada, P$_{1}$P$_{2}$ = $\surd$(a$^2$+49b$^2$)/4, comparável às dimensões da órbita, o campo da estrela ($\propto$1/r$^2$) não é uniforme na região. Nesse caso o CM não segue a órbita original. O Poliedro aponta essa inconsistência e sugere anular a questão. Se a banca anular ou aceitar outras interpretações, a nota dos dois itens pode mudar.%
\item%
\textbf{Q5} \textit{(solution)} No item b, a resolução não inclui no calor absorvido a energia para erguer o centro de massa da própria viga, mg$\Delta$L/2. O gabarito Poliedro inclui. A omissão é fisicamente incorreta, embora seja numericamente desprezível (variação relativa de cerca de 10$^{-}$$^{8}$ em $\eta$). Se a banca exigir a expressão completa, isso pode custar parte da nota do item b. A resolução também é incoerente ao manter o termo M1g$\alpha$L, que tem a mesma ordem de grandeza desprezível.%
\end{itemize}%
\section*{Questão 1}%
\href{https://poliedroresolve.sistemapoliedro.com.br/vestibulares/ita/2026/ita-2-fase-dia-3-30-10-2025-2-fase-dia-3/dissertativa-1-fisica-2-fase-dia-3-ita-2026}{Resolução do gabarito}\par\smallskip%
\begin{longtable}{@{}p{0.06\textwidth} p{0.33\textwidth} p{0.33\textwidth} p{0.1\textwidth} p{0.08\textwidth}@{}}\toprule Item & Resolução & Gabarito & Veredito & Certo \\ \midrule
a & R\_min = v$^2$(1 - $\mu$ tan$\theta$)/{[}g(tan$\theta$ + $\mu$){]} = v$^2$/{[}g tan($\theta$+$\varphi$){]}, com tan$\varphi$ = $\mu$. Vale para $\mu$ tan$\theta$ < 1; se $\mu$ tan$\theta$ $\geq$ 1, não há mínimo. No limite, N\_interna = 0 e f\_externa = $\mu$N\_externa. & r\_min = (v$^2$/g)$\cdot$(cos$\theta$ - $\mu$ sen$\theta$)/(sen$\theta$ + $\mu$ cos$\theta$). No limite, a roda interna está na iminência de perder o contato (N\_i $\approx$ 0, Fat\_i $\approx$ 0) e Fat\_E = $\mu$N\_E. & equivalente & ambos \\
b & R\_min = v$^2$/(g tan$\theta$) & r\_min = (v$^2$/g)$\cdot$cotg $\theta$ & equivalente & ambos \\
c & R\_min = v$^2$/($\mu$g) & r\_min = v$^2$/($\mu$g) & igual & ambos \\
\bottomrule\end{longtable}%
\begin{itemize}%
\item%
\textbf{a)} Dividindo numerador e denominador da expressão do gabarito por cos$\theta$, obtém{-}se (1 - $\mu$ tan$\theta$)/(tan$\theta$ + $\mu$) = 1/tan($\theta$+$\varphi$). Conferi numericamente, para vários $\theta$ e $\mu$, que as três formas coincidem. A física também é a mesma. A normal fica inclinada de $\theta$ em relação à vertical, para o centro do trem, o que bate com a figura: o raio da roda é maior junto ao eixo. O atrito na roda externa aponta para baixo, ao longo da geratriz, rumo ao centro da curva. O limite é N\_i = 0 com Fat\_E = $\mu$N\_E, e as equações de Newton são idênticas às do gabarito. A resolução ainda prova, pelo cone de atrito, que esse é de fato o máximo de v$^2$/(Rg) para qualquer distribuição de forças com N $\geq$ 0. Também acrescenta a restrição $\mu$ tan$\theta$ < 1, que o gabarito omite e que não muda a resposta.%
\item%
\textbf{b)} Como 1/tan$\theta$ = cotg $\theta$, as respostas são a mesma. A resolução deduz o resultado das equações sem atrito, (N1+N2)cos$\theta$ = Mg e (N1-N2)sen$\theta$ = Mv$^2$/R, impondo N2 $\geq$ 0. O gabarito chega a ele fazendo $\mu$ = 0 em (a). Os dois caminhos estão corretos.%
\item%
\textbf{c)} O resultado é idêntico. O gabarito faz $\theta$ = 0 em (a). A resolução faz a dedução direta: f1 + f2 $\leq$ $\mu$(N1+N2) = $\mu$Mg.%
\end{itemize}%
\paragraph{Pontos de dúvida}%
\begin{itemize}%
\item%
\textit{(statement)} No item (c), com rodas cilíndricas e o trem tratado como partícula, as normais são verticais e nenhuma roda 'perde o contato' no sentido de N $\to$ 0. O que limita R é o escorregamento lateral. As duas respostas leem 'perder o contato com o trilho' como 'sair do trilho' (escorregar ou descarrilar). Essa leitura é razoável, mas não está explícita no enunciado.%
\item%
\textit{(statement)} O enunciado não dá a altura do centro de massa nem a bitola. Por isso as duas soluções desprezam o equilíbrio de torques (tombamento) e tratam o trem como partícula. A hipótese é necessária e foi assumida pelos dois, mas o gabarito não a declara.%
\item%
\textit{(key)} O gabarito não menciona que a fórmula de (a) só tem sentido se $\mu$ tan$\theta$ < 1, ou seja, $\theta$ + arctg $\mu$ < 90$^\circ$. Se $\mu$ tan$\theta$ $\geq$ 1, o resultado seria nulo ou negativo. A resolução trata esse caso. Para valores realistas, a restrição é irrelevante.%
\item%
\textit{(key)} A transcrição da figura do gabarito diz que Fat\_E aponta 'para baixo/para fora'. Já as equações do gabarito (+Fat\_E cos$\theta$ na direção centrípeta e -Fat\_E sen$\theta$ na vertical) correspondem a um atrito para baixo e para o centro da curva, como na resolução. Isso parece descrever o sentido no desenho, não um erro físico. As equações estão corretas.%
\end{itemize}%
\section*{Questão 2}%
\href{https://poliedroresolve.sistemapoliedro.com.br/vestibulares/ita/2026/ita-2-fase-dia-3-30-10-2025-2-fase-dia-3/dissertativa-2-fisica-2-fase-dia-3-ita-2026}{Resolução do gabarito}\par\smallskip%
\begin{longtable}{@{}p{0.06\textwidth} p{0.33\textwidth} p{0.33\textwidth} p{0.1\textwidth} p{0.08\textwidth}@{}}\toprule Item & Resolução & Gabarito & Veredito & Certo \\ \midrule
a & N\_A = g cos$\theta$ {[}M/2 + (m(L-a-d) + m\_C(b-d))/(L-2d){]}; N\_B = g cos$\theta$ {[}M/2 + (m(a-d) + m\_C(L-b-d))/(L-2d){]} & N\_A = {[}m(L-a-d) + M(L/2-d) + m\_C(b-d){]}$\cdot$g cos$\theta$/(L-2d); N\_B = {[}m(a-d) + M(L/2-d) + m\_C(L-b-d){]}$\cdot$g cos$\theta$/(L-2d) & equivalente & ambos \\
b & f\_A + f\_B = (M+m+m\_C) g sen$\theta$, rampa acima. A resolução reconhece que a estática só determina a soma. Depois supõe o mesmo $\mu$ nos dois apoios e deslizamento iminente ($\mu$ = $\mu$\_min = tan$\theta$) e chega a f\_A = N\_A tan$\theta$ e f\_B = N\_B tan$\theta$, que são os valores em destaque. & Fat\_A + Fat\_B = (m+M+m\_C) g sen$\theta$, ambas paralelas ao plano e para cima. Não há dados para calcular cada atrito separadamente (problema indeterminado). & diferente & gabarito \\
\bottomrule\end{longtable}%
\begin{itemize}%
\item%
\textbf{a)} São a mesma expressão, porque M(L/2-d)/(L-2d) = M/2. Os dois tomam momentos em A e em B com os mesmos braços: (a-d), (L/2-d) e (L-b-d) em relação a A; (L-a-d), (L/2-d) e (b-d) em relação a B. Conferi numericamente com valores aleatórios. As expressões coincidem e o sistema completo ($\Sigma$Fx, $\Sigma$Fy e $\Sigma$$\tau$ em coordenadas globais) fica em equilíbrio, com N\_A + N\_B = (M+m+m\_C)g cos$\theta$.%
\item%
\textbf{b)} A soma dos atritos é igual nos dois e está correta. A diferença está na conclusão sobre cada apoio. Com espessura desprezível, os atritos e as componentes paralelas dos pesos estão sobre a mesma linha de ação e não entram na equação de momentos. Assim, qualquer par (f\_A, f\_B) com a soma certa satisfaz todas as equações de equilíbrio, e conferi isso numericamente. O problema é hiperestático na direção do plano, como diz o gabarito. A divisão f\_i = N\_i tan$\theta$ da resolução é dedução correta só sob duas hipóteses que o enunciado não dá: $\mu$ igual nos dois apoios e $\mu$ exatamente igual a tan$\theta$, isto é, deslizamento iminente. O enunciado não fornece $\mu$ e não diz que a barra está na iminência de deslizar. Para $\mu$ > tan$\theta$ a divisão volta a ser indeterminada. Outro modelo igualmente plausível, com apoios de mesma rigidez ao cisalhamento, daria f\_A = f\_B = (M+m+m\_C)g sen$\theta$/2. A resolução declara a hipótese e afirma que só a soma é determinada, então contém a resposta do gabarito. Mesmo assim, os valores em destaque não decorrem do enunciado. Deve ganhar crédito pela soma, mas os valores individuais não se sustentam como resposta geral.%
\end{itemize}%
\paragraph{Pontos de dúvida}%
\begin{itemize}%
\item%
\textit{(statement)} O item b pede para "calcular a força de atrito em cada apoio" e lista os parâmetros permitidos sem coeficiente de atrito. Isso sugere que a banca esperava valores individuais, talvez proporcionais às normais (f\_i = N\_i tan$\theta$, como na resolução) ou divididos igualmente. Pela estática de corpo rígido, porém, só a soma é determinada, como diz o Poliedro. Vale levar ao professor qual critério a banca do ITA deve usar. Se o gabarito oficial der f\_i = N\_i tan$\theta$, a resolução fica com a resposta esperada.%
\item%
\textit{(solution)} No item b, a resolução põe em destaque f\_A = N\_A tan$\theta$ e f\_B = N\_B tan$\theta$ com base na hipótese de deslizamento iminente ($\mu$ = tan$\theta$) e $\mu$ igual nos dois apoios. Nenhuma das duas vem do enunciado. A ressalva aparece no texto, mas a resposta final dá a impressão de um resultado determinado. Pelo critério do gabarito Poliedro (indeterminado), esses valores não se justificam.%
\item%
\textit{(solution)} Os dois desprezam a espessura da barra e a altura dos apoios. A figura mostra apoios triangulares com altura visível. Se a linha da barra estiver a uma altura h acima dos pontos de contato, as componentes paralelas dos pesos geram momento (M+m+m\_C)g sen$\theta$$\cdot$h em relação aos apoios, e as normais mudam. A resolução declara essa hipótese; o gabarito a usa sem dizer. É uma hipótese razoável, já que o enunciado não dá h.%
\item%
\textit{(statement)} O item a fala em "forças normais exercidas sobre os apoios". Os dois calculam a normal que o plano exerce sobre a barra nos apoios, que pela 3ª lei tem o mesmo módulo. A diferença é só de redação.%
\end{itemize}%
\section*{Questão 3}%
\href{https://poliedroresolve.sistemapoliedro.com.br/vestibulares/ita/2026/ita-2-fase-dia-3-30-10-2025-2-fase-dia-3/dissertativa-3-fisica-2-fase-dia-3-ita-2026}{Resolução do gabarito}\par\smallskip%
\begin{longtable}{@{}p{0.06\textwidth} p{0.33\textwidth} p{0.33\textwidth} p{0.1\textwidth} p{0.08\textwidth}@{}}\toprule Item & Resolução & Gabarito & Veredito & Certo \\ \midrule
a & F\_r = M{[}$\omega$$^2$(L-2a)/2 + g cos $\theta${]}, ao longo da haste e apontando para o pivô, isto é, do CM para o lado do trecho curto. Em vetor: R\_r = -M{[}$\omega$$^2$(L/2-a) + g cos $\theta${]} r?, com r? = (sen $\theta$, -cos $\theta$) indo do pivô para o CM. & P\_R = Mg cos $\theta$ + M$\omega$$^2$(L-2a)/2, ao longo da haste e apontando para a extremidade mais próxima do pivô (trecho de comprimento a). Calculado no referencial girante, com a < L/2. & igual & ambos \\
b & $\theta$ = arccos{[}$\omega$$^2$(2a-L)/(2g){]}. Existe se $\omega$$^2$(L-2a) $\leq$ 2g; para a < L/2 o ângulo fica entre 90$^\circ$ e 180$^\circ$. & $\theta$ = arccos{[}-$\omega$$^2$(L-2a)/(2g){]}, com a < L/2. & equivalente & ambos \\
\bottomrule\end{longtable}%
\begin{itemize}%
\item%
\textbf{a)} As expressões e o sentido coincidem. O sentido 'para o pivô, vindo do CM' é o mesmo que 'para o trecho curto'. Refiz a conta pelos dois caminhos. (1) Teorema do CM projetado em r?: R\_r + Mg cos $\theta$ = -M$\omega$$^2$d, com d = L/2 - a. (2) Referencial girante: $\int$(M/L)$\omega$$^2$ r dr de -a até L-a dá M$\omega$$^2$(L-2a)/2, igual à soma do gabarito, M$\omega$$^2${[}(L-a)$^2$ - a$^2${]}/(2L). Conferi numericamente (M=2, L=1,5, a=0,3, $\omega$=2,7, $\theta$=0,8): os dois dão 20,495 N. Os dois também notam que $\alpha$ não aparece na direção radial.%
\item%
\textbf{b)} $\omega$$^2$(2a-L) = -$\omega$$^2$(L-2a), então o argumento do arccos é o mesmo. Conferido numericamente: os dois dão 1,9050 rad. A resolução acrescenta a condição de existência, $\omega$ $\leq$ $\surd$(2g/(L-2a)), que o gabarito não menciona.%
\end{itemize}%
\paragraph{Pontos de dúvida}%
\begin{itemize}%
\item%
\textit{(solution)} No item a, a resolução chama F\_r = M{[}$\omega$$^2$(L-2a)/2 + g cos $\theta${]} de 'módulo, orientada para o pivô'. Isso só vale enquanto a expressão é positiva. Para $\theta$ > 90$^\circ$ com $\omega$ pequeno, ela fica negativa e a força aponta para fora do pivô. O certo é chamá{-}la de componente com sinal, positiva no sentido do pivô, que é como o gabarito trata P\_R. Não muda a expressão.%
\item%
\textit{(key)} O gabarito não diz quando o item b tem solução. O arccos só existe se $\omega$$^2$(L-2a) $\leq$ 2g. A resolução cita essa condição e o gabarito não.%
\item%
\textit{(statement)} Os dois partem de a < L/2, como na figura; o enunciado não diz isso explicitamente. Além disso, $\theta$ é um ângulo sem sinal em relação à vertical. A posição espelhada (-$\theta$, com a haste do outro lado da vertical) também anula a componente radial. O arccos devolve só o representante entre 0 e $\pi$.%
\end{itemize}%
\section*{Questão 4}%
\href{https://poliedroresolve.sistemapoliedro.com.br/vestibulares/ita/2026/ita-2-fase-dia-3-30-10-2025-2-fase-dia-3/dissertativa-4-fisica-2-fase-dia-3-ita-2026}{Resolução do gabarito}\par\smallskip%
\begin{longtable}{@{}p{0.06\textwidth} p{0.33\textwidth} p{0.33\textwidth} p{0.1\textwidth} p{0.08\textwidth}@{}}\toprule Item & Resolução & Gabarito & Veredito & Certo \\ \midrule
a & y$_{2}$ = 2b ou y$_{2}$ = -4b. Vem de m$_{1}$ = 4m/5, m$_{2}$ = m/5 e CM = (4a/5, (b+y$_{2}$)/5), com o CM sobre a elipse original. & y$_{2}$ = 2b ou y$_{2}$ = -4b, com a mesma premissa de que o CM continua na elipse original. O Poliedro diz que essa premissa é falsa e sugere anular a questão. & igual & ambos \\
b & Estrela no foco (c,0) e y$_{2}$ = 2b: U\_min = -(GMm/5){[}16/$\surd$((3a-4c)$^2$+b$^2$) + 1/$\surd$((a-c)$^2$+4b$^2$){]} - 16Gm$^2$/(25$\surd$(a$^2$+49b$^2$)), com c = $\surd$(a$^2$-b$^2$). A resolução considera o último termo desprezível para M ? m e dá também a forma aproximada só com os termos da estrela. & E\_pg,mín = -GMm$_{2}${[}4/$\surd$((3a/4-$\surd$(a$^2$-b$^2$))$^2$ + (b/4)$^2$) + 1/$\surd$((a-$\surd$(a$^2$-b$^2$))$^2$ + (2b)$^2$){]}, com y$_{2}$ = 2b e a estrela em (c,0). & equivalente & ambos \\
\bottomrule\end{longtable}%
\begin{itemize}%
\item%
\textbf{a)} Refiz a conta. x\_CM = (0,8$\cdot$3a/4 + 0,2$\cdot$a) = 4a/5. Na elipse, 16/25 + y\_CM$^2$/b$^2$ = 1, então y\_CM = $\pm$3b/5 e (b+y$_{2}$)/5 = $\pm$3b/5, o que dá y$_{2}$ = 2b ou y$_{2}$ = -4b. A resolução e o gabarito chegam ao mesmo resultado pelo mesmo caminho.%
\item%
\textbf{b)} As duas escolhem a mesma configuração: foco (c,0) e y$_{2}$ = 2b. Testei numericamente as 4 combinações (foco $\pm$c, y$_{2}$ = 2b ou -4b) em 2000 pares (a,b) sorteados, e o mínimo foi sempre esse. Os termos da estrela são idênticos: m$_{2}$ = m/5 e 4/d$_{1}$ = 16/$\surd$((3a-4c)$^2$+b$^2$), porque d$_{1}$ = $\surd$((3a-4c)$^2$+b$^2$)/4. A igualdade foi conferida numericamente. As formas alternativas dos radicandos (25a$^2$-24ac-15b$^2$ e 2a$^2$-2ac+3b$^2$) também estão certas. A única diferença é o termo de interação entre m$_{1}$ e m$_{2}$, -Gm$_{1}$m$_{2}$/d = -16Gm$^2$/(25$\surd$(a$^2$+49b$^2$)), com d = $\surd$(a$^2$+49b$^2$)/4 (conferido). A resolução inclui esse termo, que pertence de fato à energia do sistema, e o declara desprezível por ser da ordem de m/M; o gabarito o omite. A forma aproximada da resolução coincide com a do gabarito. Incluir o termo não muda a configuração ótima, porque y$_{2}$ = 2b também minimiza d.%
\end{itemize}%
\paragraph{Pontos de dúvida}%
\begin{itemize}%
\item%
\textit{(statement)} A questão depende de uma premissa que o enunciado não declara: o CM de m$_{1}$+m$_{2}$ continua sobre a elipse original. Com a separação dada, P$_{1}$P$_{2}$ = $\surd$(a$^2$+49b$^2$)/4, comparável às dimensões da órbita, o campo da estrela ($\propto$1/r$^2$) não é uniforme na região. Nesse caso o CM não segue a órbita original. O Poliedro aponta essa inconsistência e sugere anular a questão. Se a banca anular ou aceitar outras interpretações, a nota dos dois itens pode mudar.%
\item%
\textit{(solution)} A resolução justifica a premissa dizendo que as forças internas 'não alteram o movimento do CM' e a chama de 'idealização habitual'. Ela até reconhece que isso equivale a desprezar a diferença entre as forças da estrela sobre as partes e a força sobre m concentrada no CM. Mas não observa que, com a separação dada no enunciado, essa aproximação deixa de valer, como faz o Poliedro. Isso não altera a resposta pretendida, mas a justificativa física é frágil.%
\item%
\textit{(key)} O gabarito calcula a 'energia potencial gravitacional do sistema' só com as interações estrela--m$_{1}$ e estrela--m$_{2}$ e omite, sem comentar, o termo -Gm$_{1}$m$_{2}$/d = -16Gm$^2$/(25$\surd$(a$^2$+49b$^2$)). A resolução inclui esse termo e o declara desprezível para M ? m, o que é mais completo. Fica em aberto se a banca espera o termo de interação ou não.%
\item%
\textit{(statement)} O enunciado não diz em qual foco está a estrela nem qual dos dois valores de y$_{2}$ vale. Por isso o item b pede a 'menor energia possível'. Tanto a resolução quanto o gabarito resolvem essa ambiguidade do mesmo modo: foco (c,0) e y$_{2}$ = 2b.%
\end{itemize}%
\section*{Questão 5}%
\href{https://poliedroresolve.sistemapoliedro.com.br/vestibulares/ita/2026/ita-2-fase-dia-3-30-10-2025-2-fase-dia-3/dissertativa-5-fisica-2-fase-dia-3-ita-2026}{Resolução do gabarito}\par\smallskip%
\begin{longtable}{@{}p{0.06\textwidth} p{0.33\textwidth} p{0.33\textwidth} p{0.1\textwidth} p{0.08\textwidth}@{}}\toprule Item & Resolução & Gabarito & Veredito & Certo \\ \midrule
a & W = (M1 - M2)$\cdot$g$\cdot$$\alpha$$\cdot$L$\cdot$$\Delta$T = 700$\cdot$$\alpha$$\cdot$L$\cdot$(M1 - M2) J (SI), com $\Delta$T = 70 K e g = 10 m/s$^2$ & W = (M1 - M2)$\cdot$g$\cdot$$\alpha$$\cdot$L$\cdot$(90 - 20) = 70$\cdot$(M1 - M2)$\cdot$g$\cdot$$\alpha$$\cdot$L & equivalente & ambos \\
b & $\eta$ = (M1 - M2)g$\alpha$L / ($\rho$Vc + 3Rm/M\_Cu + M1g$\alpha$L) $\approx$ 10$\alpha$L(M1 - M2)/($\rho$Vc + 3Rm/M\_Cu) & $\eta$ = (M1 - M2)g$\alpha$L / ($\rho$Vc + 3Rm/M\_Cu + M1g$\alpha$L + mg$\alpha$L/2) & diferente & gabarito \\
\bottomrule\end{longtable}%
\begin{itemize}%
\item%
\textbf{a)} Mesmo raciocínio: M1 sobe $\Delta$L = $\alpha$L$\Delta$T no aquecimento e M2 desce $\Delta$L no resfriamento, então o trabalho líquido é (M1 - M2)g$\Delta$L. A resolução só substituiu g = 10 m/s$^2$ (70$\cdot$10 = 700), o que é legítimo porque g foi dado na prova. As duas formas são equivalentes.%
\item%
\textbf{b)} A diferença está só no denominador. No aquecimento, a própria viga se dilata e seu centro de massa sobe $\Delta$L/2. Essa energia potencial, mg$\Delta$L/2 = mg$\alpha$L$\Delta$T/2, também sai do calor fornecido pela fonte. A resolução contou apenas a elevação de M1 e deixou de fora esse termo. Pela 1ª lei, o gabarito está mais completo: Q\_q = $\Delta$T($\rho$Vc + 3Rm/M\_Cu + M1g$\alpha$L + mg$\alpha$L/2). No ciclo, esse termo se anula, pois o centro de massa da viga desce $\Delta$L/2 no resfriamento. Por isso o trabalho líquido do item a não muda, e a verificação Q\_abs - Q\_ced = W da resolução continua valendo. Na prática, a diferença é desprezível. No exemplo da própria resolução (viga de 1 kg, L = 1 m, $\alpha$ = 1,7$\times$10$^{-}$$^{5}$ K$^{-}$$^{1}$), o termo omitido vale cerca de 8,5$\times$10$^{-}$$^{5}$ J/K, contra cerca de 4,6$\times$10$^3$ J/K de capacidade térmica, e as duas eficiências diferem relativamente em cerca de 2$\times$10$^{-}$$^{8}$ (conferido em python). Se todos os termos mecânicos forem desprezados, a forma aproximada da resolução coincide com o limite do gabarito. Ainda assim, a expressão 'exata' da resolução é incompleta e incoerente: ela mantém M1g$\alpha$L, que tem a mesma ordem de grandeza desprezível, mas omite mg$\alpha$L/2.%
\end{itemize}%
\paragraph{Pontos de dúvida}%
\begin{itemize}%
\item%
\textit{(solution)} No item b, a resolução não inclui no calor absorvido a energia para erguer o centro de massa da própria viga, mg$\Delta$L/2. O gabarito Poliedro inclui. A omissão é fisicamente incorreta, embora seja numericamente desprezível (variação relativa de cerca de 10$^{-}$$^{8}$ em $\eta$). Se a banca exigir a expressão completa, isso pode custar parte da nota do item b. A resolução também é incoerente ao manter o termo M1g$\alpha$L, que tem a mesma ordem de grandeza desprezível.%
\item%
\textit{(key)} O gabarito, assim como a resolução, despreza a dilatação da água: o trabalho contra a atmosfera e a elevação do centro de massa da água. Também despreza a capacidade térmica do recipiente. São hipóteses razoáveis que não foram declaradas. Dentro da mesma lógica que inclui mg$\Delta$L/2, a variação de energia potencial da água poderia ser tão relevante quanto esse termo.%
\item%
\textit{(statement)} O enunciado não diz se a eficiência deve considerar a forma exata ou se os termos mecânicos podem ser desprezados diante das capacidades térmicas. A forma aproximada da resolução, $\eta$ $\approx$ g$\alpha$L(M1 - M2)/($\rho$Vc + 3Rm/M\_Cu), provavelmente seria aceita, mas isso depende do critério da banca.%
\item%
\textit{(solution)} A resolução usa c\_P $\approx$ c\_V para o sólido e trata os processos como quase estáticos, com a viga sempre em equilíbrio térmico com o banho. O gabarito adota as mesmas hipóteses, sem declará{-}las.%
\end{itemize}%
\section*{Questão 6}%
\href{https://poliedroresolve.sistemapoliedro.com.br/vestibulares/ita/2026/ita-2-fase-dia-3-30-10-2025-2-fase-dia-3/dissertativa-6-fisica-2-fase-dia-3-ita-2026}{Resolução do gabarito}\par\smallskip%
\begin{longtable}{@{}p{0.06\textwidth} p{0.33\textwidth} p{0.33\textwidth} p{0.1\textwidth} p{0.08\textwidth}@{}}\toprule Item & Resolução & Gabarito & Veredito & Certo \\ \midrule
a & L = f1 + f2 = 85 cm (ajuste normal, sistema afocal) & L = f1 + f2 = 80 + 5 = 85 cm & igual & ambos \\
b & G = -f1/f2 = -16 (módulo 16); imagem invertida (virtual, no infinito) & G = 16 em módulo; com sinal, G = -16; imagem invertida & igual & ambos \\
c & p'1 = -80 cm (virtual, direita, A1 = +2); p2 = 165 cm; p'2 = 165/32 $\approx$ 5,2 cm depois da ocular; A = -1/16; imagem final real, invertida e menor (1/16 do objeto) & p'1 = -80 cm; p2 = 165 cm; p'2 $\approx$ 5,2 cm; A $\approx$ -0,063; imagem real, invertida, menor, a 5,2 cm da ocular & igual & ambos \\
\bottomrule\end{longtable}%
\begin{itemize}%
\item%
\textbf{a)} Para que o objeto e a imagem final fiquem no infinito, o foco imagem da objetiva tem de coincidir com o foco objeto da ocular. Daí L = 80 + 5 = 85 cm.%
\item%
\textbf{b)} Os dois chegam a |G| = f1/f2 = 80/5 = 16, com sinal negativo, e dizem que a imagem é invertida. A resolução tira o resultado da construção dos raios (y1 = -f1$\cdot$$\theta$, $\theta$' = y1/f2). O gabarito usa a fórmula geral G = (f1/f2)(1 + f2/D) e faz D $\to$ $\infty$.%
\item%
\textbf{c)} Refiz as contas com frações exatas: p'1 = -80 cm, A1 = +2, p2 = 85 + 80 = 165 cm, p'2 = 165/32 = 5,15625 cm, A2 = -1/32 e A = -1/16 = -0,0625. O -0,063 do gabarito é só esse valor arredondado. A resolução ainda confere o resultado de forma independente: num sistema afocal o aumento lateral vale -f2/f1 = -1/16 para qualquer posição do objeto.%
\end{itemize}%
\paragraph{Pontos de dúvida}%
\begin{itemize}%
\item%
\textit{(statement)} O item (c) não diz se o tubo continua com L = 85 cm, o ajuste do item (a). A resolução e o gabarito supõem que sim, o que é a leitura natural. Se a ocular fosse reajustada para focalizar o objeto próximo, a resposta mudaria.%
\item%
\textit{(key)} O gabarito arredonda p'2 para 5,2 cm antes de calcular A, e chega a -0,063. O valor exato é -1/16 = -0,0625. A diferença é pequena e não muda nenhuma conclusão.%
\item%
\textit{(statement)} Convenção de sinal do aumento angular. Muitos livros definem G em módulo (G = 16) e indicam a inversão à parte. Como os dois dão o módulo e também o sinal (-16), isso não afeta a correção.%
\end{itemize}%
\section*{Questão 7}%
\href{https://poliedroresolve.sistemapoliedro.com.br/vestibulares/ita/2026/ita-2-fase-dia-3-30-10-2025-2-fase-dia-3/dissertativa-7-fisica-2-fase-dia-3-ita-2026}{Resolução do gabarito}\par\smallskip%
\begin{longtable}{@{}p{0.06\textwidth} p{0.33\textwidth} p{0.33\textwidth} p{0.1\textwidth} p{0.08\textwidth}@{}}\toprule Item & Resolução & Gabarito & Veredito & Certo \\ \midrule
a & I = I0$\cdot$x\^{}(N+1)$\cdot$cos\^{}(2N)($\varphi$/N) & I\_final = I0$\cdot$x\^{}(N+1)$\cdot$cos\^{}(2N)(?/N) & igual & ambos \\
b & r = sqrt( P / ($\pi$$\cdot$I0$\cdot$x\^{}(N+1)$\cdot$cos\^{}(2N)($\varphi$/N)) ) & r = sqrt( P / ($\pi$$\cdot$I0$\cdot$x\^{}(N+1)$\cdot$cos\^{}(2N)(?/N)) ) & igual & ambos \\
\bottomrule\end{longtable}%
\begin{itemize}%
\item%
\textbf{a)} As duas respostas são idênticas. Refiz a conta. No 1º polarizador o ângulo é 0, então a lei de Malus deixa passar tudo e só entra a perda: I1 = x$\cdot$I0. Em cada um dos N polarizadores seguintes, a luz chega girada de $\beta$ em relação ao eixo, o que dá I\_\{k+1\} = x$\cdot$cos$^2$$\beta$$\cdot$I\_k. Assim, I = x$\cdot$I0$\cdot$(x cos$^2$$\beta$)\^{}N = I0$\cdot$x\^{}(N+1)$\cdot$cos\^{}(2N)($\varphi$/N). A transcrição confere com o original (p. 8).%
\item%
\textbf{b)} As duas respostas são idênticas. Como o feixe é cilíndrico e homogêneo, P = I$\cdot$$\pi$r$^2$, logo r = sqrt(P/($\pi$I)), com o I do item a. As unidades conferem (m). A resolução ainda observa, corretamente, que a fórmula exige cos($\varphi$/N) $\neq$ 0.%
\end{itemize}%
\paragraph{Pontos de dúvida}%
\begin{itemize}%
\item%
\textit{(statement)} A frase "sua intensidade após a passagem por cada polarizador é igual a uma fração x daquela que seria obtida em um caso ideal" tem uma segunda leitura possível: x aplicado uma única vez, sobre a intensidade do sistema todo ideal, o que daria I = x$\cdot$I0$\cdot$cos\^{}(2N)($\varphi$/N). A resolução e o gabarito usam a leitura mais física: perda x em cada polarizador, acumulada, o que dá x\^{}(N+1). Essa deve ser a resposta esperada.%
\item%
\textit{(statement)} O "analisador" foi tratado como um medidor de potência, sem efeito sobre o feixe. Se fosse mais um polarizador, alinhado ao último e também com perda, entraria um fator x a mais. Gabarito e resolução fazem a mesma hipótese, que é a natural.%
\item%
\textit{(solution)} A redação "Como I $\leq$ I0, a energia se conserva" é imprecisa. I $\leq$ I0 só mostra que o resultado é compatível com a conservação de energia (não há ganho). A energia luminosa não se conserva: parte é absorvida ou refletida. Isso não afeta as respostas.%
\end{itemize}%
\section*{Questão 8}%
\href{https://poliedroresolve.sistemapoliedro.com.br/vestibulares/ita/2026/ita-2-fase-dia-3-30-10-2025-2-fase-dia-3/dissertativa-8-fisica-2-fase-dia-3-ita-2026}{Resolução do gabarito}\par\smallskip%
\begin{longtable}{@{}p{0.06\textwidth} p{0.33\textwidth} p{0.33\textwidth} p{0.1\textwidth} p{0.08\textwidth}@{}}\toprule Item & Resolução & Gabarito & Veredito & Certo \\ \midrule
a & P\_min = (36000 J + 6000 J)/1200 s = 35 W & P\_min = 42000 J / 1200 s = 35 W & igual & ambos \\
b & $\Delta$t\_max = L$\cdot$$\Delta$I/$\varepsilon$ = 8$\pi$$^2$$\cdot$10$^{-}$$^{7}$ $\cdot$ 14/5 = 2,24$\pi$$^2$$\cdot$10$^{-}$$^{6}$ s $\approx$ 2,2$\cdot$10$^{-}$$^{5}$ s & $\Delta$t\_max = (56/25)$\cdot$$\pi$$^2$$\cdot$10$^{-}$$^{6}$ s $\approx$ 2,2$\cdot$10$^{-}$$^{5}$ s & igual & ambos \\
\bottomrule\end{longtable}%
\begin{itemize}%
\item%
\textbf{a)} O enunciado transcrito confere com a página 9 do PDF. As duas resoluções seguem o mesmo raciocínio. A bateria armazena q$\varepsilon$ = 7200 C $\cdot$ 5,0 V = 36000 J e o aparelho absorve mc$\Delta$$\theta$ = 100 $\cdot$ 2,0 $\cdot$ 30 = 6000 J. A soma, 42000 J, dividida por 1200 s dá 35 W, conferido em Python.%
\item%
\textbf{b)} 56/25 = 2,24, então os dois valores são iguais. Refiz a conta: L = $\mu$0$\cdot$n$^2$$\cdot$$\pi$$\cdot$r$^2$$\cdot$? = 4$\pi$$\cdot$10$^{-}$$^{7}$ $\cdot$ (2000)$^2$ $\cdot$ $\pi$ $\cdot$ (5$\cdot$10$^{-}$$^3$)$^2$ $\cdot$ 2$\cdot$10$^{-}$$^2$ = 8$\pi$$^2$$\cdot$10$^{-}$$^{7}$ H $\approx$ 7,9$\cdot$10$^{-}$$^{6}$ H. Daí $\Delta$t = L$\cdot$14/5 $\approx$ 2,211$\cdot$10$^{-}$$^{5}$ s. A expressão do gabarito, N$\cdot$$\mu$0$\cdot$n$\cdot$$\pi$$\cdot$r$^2$$\cdot$$\Delta$I/$\varepsilon$ com N = 40, é a mesma coisa, porque N$\cdot$n = n$^2$$\cdot$?. Muda só a justificativa do 'máximo'. A resolução usa M $\leq$ $\surd$(L1$\cdot$L2) = L, com igualdade no acoplamento perfeito, que é a justificativa mais rigorosa. O gabarito argumenta que um $\Delta$t maior induziria menos de 5 V. A resolução ainda faz uma verificação: 5 V $\times$ 7 A (corrente média na rampa) = 35 W, coerente com o item a.%
\end{itemize}%
\paragraph{Pontos de dúvida}%
\begin{itemize}%
\item%
\textit{(key)} O gabarito justifica o 'máximo' dizendo que um $\Delta$t maior induziria menos de 5 V. No modelo ideal, porém, a condição $\varepsilon$ = M$\cdot$$\Delta$I/$\Delta$t determina $\Delta$t de forma única. O caráter de máximo vem de M $\leq$ $\surd$(L1$\cdot$L2) = L, ou de o modelo de solenoide ideal superestimar o campo. A resolução apresenta essa justificativa e o gabarito não. A justificativa do gabarito é fraca, mas não muda o valor.%
\item%
\textit{(solution)} Na verificação energética, a resolução afirma que 'cada ciclo transfere no máximo $\frac12$$\cdot$L$\cdot$I$^2$\_max'. Isso é discutível: num transformador com carga, a corrente no secundário altera o fluxo e a energia transferida não se limita à energia magnética armazenada. Além disso, a conta supõe ciclos em dente de serra sem intervalo entre eles. Como o número bate ($\frac12$$\cdot$L$\cdot$I$^2$/($\varepsilon$$\cdot$I/2) = L$\cdot$I/$\varepsilon$), é só uma verificação complementar e não afeta a resposta.%
\item%
\textit{(statement)} O solenoide é curto (?/r = 4), então B = $\mu$0$\cdot$n$\cdot$I e L = $\mu$0$\cdot$n$^2$$\cdot$$\pi$$\cdot$r$^2$$\cdot$? são aproximações. O enunciado também ignora o efeito da corrente do secundário sobre o fluxo e supõe tensão constante de 5,0 V durante toda a carga (E = q$\cdot$V). Essas hipóteses de idealização ficam implícitas tanto na resolução quanto no gabarito.%
\item%
\textit{(statement)} A prova traz $\mu$0 = 4$\pi$$\cdot$10$^{-}$$^{7}$ com unidade N$\cdot$A$^{-}$$^{1}$. O correto seria N/A$^2$ (ou H/m). É erro de impressão sem efeito no valor numérico.%
\end{itemize}%
\section*{Questão 9}%
\href{https://poliedroresolve.sistemapoliedro.com.br/vestibulares/ita/2026/ita-2-fase-dia-3-30-10-2025-2-fase-dia-3/dissertativa-9-fisica-2-fase-dia-3-ita-2026}{Resolução do gabarito}\par\smallskip%
\begin{longtable}{@{}p{0.06\textwidth} p{0.33\textwidth} p{0.33\textwidth} p{0.1\textwidth} p{0.08\textwidth}@{}}\toprule Item & Resolução & Gabarito & Veredito & Certo \\ \midrule
a & d = 600$\surd$2 nm $\approx$ 849 nm $\approx$ 8,5$\times$10$^{-}$$^{7}$ m (tg$\theta$ = 20/20 $\Rightarrow$ $\theta$ = 45$^\circ$, d$\cdot$sen45$^\circ$ = 600 nm) & d = 600$\surd$2 nm $\approx$ 848,5 nm & igual & ambos \\
b & x\_V = 20$\surd$(2/7) cm = 20$\surd$14/7 cm $\approx$ 10,7 cm (sen$\theta$ = $\surd$2/3, $\theta$ $\approx$ 28$^\circ$) & y' = 20$\surd$14/7 cm $\approx$ 10,7 cm (sen$\theta$' = $\surd$2/3) & igual & ambos \\
\bottomrule\end{longtable}%
\begin{itemize}%
\item%
\textbf{a)} Mesma física e mesma expressão exata. Conferi: 600$\surd$2 = 848,53 nm, ou seja, cerca de 1,18$\times$10$^3$ ranhuras/mm. Os valores 849 nm e 8,5$\times$10$^{-}$$^{7}$ m são só arredondamentos de 848,5 nm.%
\item%
\textbf{b)} Resultado idêntico. Conferi: sen$\theta$' = 400/(600$\surd$2) = $\surd$2/3 $\approx$ 0,4714, $\theta$' $\approx$ 28,1$^\circ$, tg$\theta$' = $\surd$(2/7) e y' = 20$\surd$(2/7) = 20$\surd$14/7 $\approx$ 10,69 cm. O gabarito chega lá por sen $\theta$' = y'/$\surd$(D$^2$+y'$^2$) e a resolução por tg$\theta$' = y'/D, que é a mesma coisa. A resolução também mostra, com razão, que a aproximação de ângulos pequenos (que daria 13,3 cm) não vale com $\theta$ = 45$^\circ$.%
\end{itemize}%
\paragraph{Pontos de dúvida}%
\begin{itemize}%
\item%
\textit{(statement)} Entende{-}se 'ponto central' como o pé da perpendicular da grade até o solo (onde fica o máximo central, m = 0), e a distância de 20 cm como medida no solo a partir desse ponto. A resolução e o gabarito leem assim, e é a única leitura natural do enunciado.%
\item%
\textit{(solution)} Hipóteses não ditas no enunciado, mas razoáveis: regime de Fraunhofer (d ? H) e Sol como fonte de ondas planas (o diâmetro angular do Sol, cerca de 0,5$^\circ$, é desprezado). Não alteram o resultado esperado.%
\end{itemize}%
\section*{Questão 10}%
\href{https://poliedroresolve.sistemapoliedro.com.br/vestibulares/ita/2026/ita-2-fase-dia-3-30-10-2025-2-fase-dia-3/dissertativa-10-fisica-2-fase-dia-3-ita-2026}{Resolução do gabarito}\par\smallskip%
\begin{longtable}{@{}p{0.06\textwidth} p{0.33\textwidth} p{0.33\textwidth} p{0.1\textwidth} p{0.08\textwidth}@{}}\toprule Item & Resolução & Gabarito & Veredito & Certo \\ \midrule
a & $\Delta$t\_Terra = 387/32 anos = 12,09375 anos $\approx$ 12,1 anos & t\_Terra = 10,75 + 1,34375 = 12,09375 anos $\approx$ 12,1 anos & igual & ambos \\
b & $\Delta$$\tau$\_2 = 387/40 anos = 9,675 anos $\approx$ 9,7 anos & t\_nave2 = 12,09375/(5/4) = 9,675 anos (tempo próprio, $\gamma$2 = 5/4) & igual & ambos \\
\bottomrule\end{longtable}%
\begin{itemize}%
\item%
\textbf{a)} Refiz a conta. A nave 1 chega a Sirius em t\_E = 8,6/0,8 = 10,75 anos. Nesse instante a nave 2 está em 6,45 a.l., portanto a 2,15 a.l. do ponto de emissão. O sinal e a nave se aproximam a 1,6c, o que leva 1,34375 ano. Total: 12,09375 = 387/32 anos (conferido em python com frações). Resolução e gabarito seguem o mesmo raciocínio e chegam ao mesmo valor.%
\item%
\textbf{b)} Os dois eventos (partida em t0 e recepção do sinal) ocorrem na própria nave 2, que tem velocidade constante. Então o intervalo é tempo próprio: $\Delta$$\tau$ = 12,09375 $\times$ 0,8 = 9,675 anos. A resolução confere esse valor pela transformação de Lorentz no referencial da nave 2: t'\_E = 6,9875 anos e x'\_E = 2,6875 a.l., que somados dão 9,675 anos. Também verifiquei em python. O resultado coincide com o do gabarito.%
\end{itemize}%
\paragraph{Pontos de dúvida}%
\begin{itemize}%
\item%
\textit{(statement)} A expressão "tempo transcorrido na Terra" foi lida como intervalo de coordenada no referencial da Terra/Sol, sem somar o tempo que a luz levaria para trazer a notícia da recepção até a Terra. A resolução deixa isso explícito e o gabarito usa a mesma leitura. É a interpretação usual, mas o enunciado não diz isso de forma explícita.%
\item%
\textit{(statement)} Hipóteses que nem o enunciado nem o gabarito declaram, e que a resolução adota: os 8,6 a.l. são medidos a partir da fronteira do sistema Solar; Sirius está em repouso no referencial da Terra; a separação transversal entre as trajetórias paralelas é desprezível; o sinal é eletromagnético e viaja a c. Nenhuma delas altera o resultado esperado.%
\end{itemize}%
\end{document}